% Page setup
 %The following defines an A4 page with classical margin distribution according to the French canon des ateliers: https://tex.stackexchange.com/questions/249190/recommended-margin-and-binding-offset
\usepackage[paperheight=297mm,paperwidth=210mm, bindingoffset=0mm,verbose,marginratio={4:6,5:7},textwidth=140mm,height=213mm,twoside, nomarginpar]{geometry}

% Text and font-related packages
\usepackage[utf8]{inputenc}
\usepackage[german, english]{babel}
\usepackage{newtxtext, newtxmath, cabin}
\usepackage[scaled=0.8]{beramono}
\usepackage[T1]{fontenc}
\usepackage{anyfontsize}

% Math, chemistry, table and figure-related packages
\usepackage{bm}
\usepackage{braket}
\usepackage{dsfont}
\usepackage{etaremune}
\usepackage{graphicx}
\usepackage{mathtools}
\usepackage{microtype}
\usepackage{multicol}
\usepackage{slantsc}
\usepackage[table]{xcolor}
\usepackage{xstring}
\usepackage[version=4]{mhchem}
\usepackage{booktabs}
\usepackage{multirow}
\usepackage{setspace}
\usepackage{lipsum}
\usepackage{enumitem}
\usepackage{pdflscape}
\usepackage{tocbasic}
\usepackage{placeins}

% Table of content style
\DeclareTOCStyleEntry[beforeskip=6pt plus 1pt minus 1pt]{tocline}{chapter}

\makeatletter

% CAPTIONS
\usepackage[font=small, labelfont=bf]{caption, subcaption}

\def\sublabels#1{\foreach \sublabel in {a,b,...,j}
    {\phantomsubcaption\label{#1_\sublabel}}}

% GRAPHICS:
\usepackage{tikz}

\usetikzlibrary{positioning}


% CHAPTER AND SECTION HEADINGS:
\usepackage[immediate]{silence}
\WarningFilter[temp]{latex}{Command} % silence the warning
\usepackage{sectsty}
\DeactivateWarningFilters[temp] % So nothing unrelated gets silenced
\allsectionsfont{\sffamily}

% EPIGRAPHS
\usepackage{epigraph}
\setlength{\epigraphwidth}{0.42\textwidth}
\def\@epitext#1{%
    \begin{minipage}{\epigraphwidth}
        \begin{\textflush}
            \itshape #1
        \end{\textflush}
        \vspace*{1ex}
    \end{minipage}
    }

\def\@episource#1{%
    \begin{minipage}{\epigraphwidth}
        \begin{\sourceflush}
            ---#1
        \end{\sourceflush}
    \end{minipage}
    }

% HEADERS AND FOOTERS
\usepackage{fancyhdr}
\usepackage[breakall, fit]{truncate}

\pagestyle{fancy}

\fancyhf{}

\def\headerline#1{\MakeUppercase{\truncate{0.7\linewidth}{#1}}}

\def\chaptermark#1{\markboth
    {\if@mainmatter\thechapter\quad\fi\headerline{#1}}
    {\headerline{#1}\if@mainmatter\quad\thechapter\fi}}

\def\sectionmark#1{\markright
    {\headerline{#1}\if@mainmatter\quad\thesection\fi}}

\fancyhead[OL]\leftmark
\fancyhead[ER]\rightmark
\fancyhead[OR]\thepage
\fancyhead[EL]\thepage

\setlength\headheight{14pt}
\def\headrulewidth{0.4pt}
\def\footrulewidth{0pt}

% COLORS
\definecolor{linkcolor}{rgb}{0.5, 0.0, 0.0}
\definecolor{basiccolor}{rgb}{0.1, 0.3, 0.5}
\definecolor{commentcolor}{rgb}{0.5, 0.5, 0.5}
\definecolor{identifiercolor}{rgb}{0.2, 0.2, 0.2}
\definecolor{keywordcolor}{rgb}{0.0, 0.4, 0.8}
\definecolor{numbercolor}{rgb}{0.5, 0.5, 0.5}
\definecolor{stringcolor}{rgb}{0.0, 0.4, 0.8}

% SOURCE CODE LISTINGS 
\usepackage{listings}
\lstdefinelanguage{f90}{
    keywords = {
        allocate,
        enddo,
        read,
        external,
        call,
        size,
        if,
        then,
        open,
        write,
        close,
        endif,
        module,
        use,
        only,
        save,
        integer,
        allocatable,
        real,
        logical,
        end,
        complex,
        do,
        cycle,
        character,
        namelist,
        return,
        subroutine,
        intent,
        },
    comment = [l]{!},
    string = [d]',
    morestring = [d]",
    sensitive = false,
    %
    escapeinside = ||,
    moredelim=**[is][\color{red}\underbar]{>>}{<<},
    }

\lstdefinelanguage{py}{
    keywords = {
        and,
        as,
        assert,
        break,
        class,
        continue,
        def,
        del,
        elif,
        else,
        except,
        exec,
        finally,
        for,
        from,
        global,
        if,
        import,
        in,
        is,
        lambda,
        not,
        or,
        pass,
        print,
        raise,
        return,
        try,
        while,
        with,
        yield,
        },
    comment = [l]{\#},
    string = [b]',
    morestring = [s]{r'}{'},
    morestring = [b]",
    morestring = [s]{r"}{"},
    sensitive = true,
    }

\lstdefinelanguage{bash}{
    keywords = {
        in,
        do,
        done,
        for,
        mpirun,
        python3,
        tee,
        test,
        },
    alsoletter = {., -, /, :},
    comment = [l]{\#},
    string = [b]',
    morestring = [b]",
    sensitive = true,
    }

\lstdefinelanguage{qe}{
    comment = [l]{!},
    string = [d]',
    }

\lstset{
    basicstyle = \color{basiccolor} \ttfamily,
    commentstyle = \color{commentcolor},
    identifierstyle = \color{identifiercolor},
    keywordstyle = \color{keywordcolor} \bfseries,
    numberstyle = \color{numbercolor} \scriptsize,
    stringstyle = \color{stringcolor},
    %
    numbers = left,
    showstringspaces = false,
    breaklines = false,
    breakatwhitespace = true,
    showlines = true,
    %
    aboveskip = 3mm,
    belowskip = 3mm,
    numbersep = 3mm,
    }

% Marco Daniel @ https://tex.stackexchange.com/questions/34094/suppress-line-numbering-for-specific-lines-in-listings-package

\def\lstnew{%
    \lst@AddToHook{OnNewLine}{%
    \def\thelstnumber{+}%
    \color{red}%
    \definecolor{commentcolor}{named}{red}%
    \definecolor{identifiercolor}{named}{red}%
    \definecolor{keywordcolor}{named}{red}%
    \definecolor{numbercolor}{named}{red}%
    \definecolor{stringcolor}{named}{red}%
    \advance\c@lstnumber-\@ne}}

\def\lstold{\expandafter\gdef\csname\@lst hk@OnNewLine\endcsname{}}

\newcommand\code[2][bash]{\text{\lstinline[language=#1]{#2}}}

% ACRONYMS:
\usepackage{acronym} % printonlyused, withpage
\usepackage{letltxmacro}
\renewcommand*{\aclabelfont}[1]{\textbf{\boldmath\acsfont{#1}}} %https://tex.stackexchange.com/questions/229930/bolding-math-acronyms


\def\acromod#1{#1}

\foreach \acromacro in {acro,acroplural} {
    \expandafter\xdef\csname\acromacro new\endcsname##1{%
        \noexpand\@ifnextchar[
            {\csname\acromacro newopt\endcsname{##1}}
            {\csname\acromacro newnop\endcsname{##1}}}
    \expandafter\xdef\csname\acromacro newopt\endcsname##1[##2]##3{%
        \csname\acromacro old\endcsname{##1}[##2]{%
            \noexpand\protect\noexpand\acromod ##3}}
    \expandafter\xdef\csname\acromacro newnop\endcsname##1##2{%
        \csname\acromacro old\endcsname{##1}{%
            \noexpand\protect\noexpand\acromod ##2}}
    }

\foreach \acrosuffix in {,f,l,s,p,fp,lp,sp,fi,lu} {
    \edef\acromacro{ac\acrosuffix}
    \edef\Acromacro{Ac\acrosuffix}
    \protected\expandafter\xdef\csname \Acromacro\endcsname{%
        \noexpand\@ifstar
            {\csname @\Acromacro star\endcsname}
            {\csname @\Acromacro nost\endcsname}}
    \expandafter\xdef\csname @\Acromacro star\endcsname##1{%
        {\noexpand\let\noexpand\acromod\noexpand\MakeUppercase%
            \csname \acromacro\endcsname*{##1}}}
    \expandafter\xdef\csname @\Acromacro nost\endcsname##1{%
        {\noexpand\let\noexpand\acromod\noexpand\MakeUppercase%
            \csname \acromacro\endcsname{##1}}}
    }

\def\ainclude#1{\include{#1}\input{#1.uac}}
\def\tinclude#1{\include{#1}\input{#1.uac}}

% DEDICATION page
\newenvironment{dedication}
  {
   \thispagestyle{empty}
   \vspace*{\stretch{1}}
   \itshape             
   \raggedleft
  }
  {\par
   \vspace{\stretch{3}}
   \clearpage          
  }


% BIBLIOGRAPHY: This uses APSREV 4.1 standards, change as needed
\usepackage[numbers, sort&compress]{natbib}
\bibliographystyle{apsrev4-1}

\def\ccite#1{Ref\IfSubStr{#1}{,}{s}{}.~\citenum{#1}}
\def\cfcite#1{\cf{\ccite{#1}}}

% HYPERLINKS (shown in red color if print=true, otherwise they stay black)
\ifprint
\usepackage[pdfa, hidelinks, allcolors=linkcolor, backref=page, pdfencoding=auto]{hyperref}
\else
\usepackage[pdfa, colorlinks, allcolors=linkcolor, backref=page, pdfencoding=auto]{hyperref}
\fi

\def\backref#1{}
\def\backrefalt#1#2#3#4{\ifcase#1 Not cited\or Cited on page~#2\else Cited on pages~#2\fi.}

\backrefparscanfalse
\def\BibitemShut#1{\csname bibitem#1\endcsname\backrefprint}%
\usepackage{hypernat}

% REFERENCES
\usepackage[capitalise]{cleveref}

\crefname{appendix}{App.}{Apps.}
\crefname{chapter}{Chap.}{Chaps.}
\crefname{equation}{Eq.}{Eqs.}
\crefname{figure}{Fig.}{Figs.}
\crefname{section}{Sec.}{Secs.}
\crefname{table}{Tab.}{Tabs.}
\Crefname{appendix}{Appendix}{Appendices}
\Crefname{chapter}{Chapter}{Chapters}
\Crefname{equation}{Equation}{Equations}
\Crefname{figure}{Figure}{Figures}
\Crefname{section}{Section}{Sections}
\Crefname{table}{Table}{Tables}

\creflabelformat{equation}{#2\textup{#1}#3}

\def\creflastconjunction{, and~}

\makeatother
\def\see#1{(#1)}
\def\cf#1{\see{cf.\@ #1}}
\def\cfref#1{\cf{\cref{#1}}}
\def\seeref#1{\see{\cref{#1}}}
\makeatletter

% Heiko Oberdiek @ https://tex.stackexchange.com/questions/340492/can-cref-be-made-to-ignore-spaces
\let\org@@cref\@cref
\renewcommand*{\@cref}[2]{%
  \begingroup
    \toks@={}%
    \comma@parse{#2}\add@cref@item
  \edef\process@me{\endgroup
    \noexpand\org@@cref{#1}{\the\toks@}%
  }\process@me
}
\newcommand*{\add@cref@item}[1]{%
  \expandafter\ifx\expandafter\\\the\toks@\\%
    \toks@{#1}%
  \else
    \toks@\expandafter{\the\toks@,#1}%
  \fi
}

\makeatother
